% chapters/07/rangkuman.tex
\bagian{Rangkuman}

\begin{rangkumanbox}
\begin{itemize}
  \item Kelas turunan \code{pygame.sprite.Sprite} menyatukan gambar, persegi
        pembatas, dan perilaku sebuah benda. Dua anggotanya wajib bernama
        \code{image} dan \code{rect}.
  \item Wadah \code{pygame.sprite.Group} menyimpan sekumpulan benda dan
        mengerjakan \code{update} serta \code{draw} untuk seluruhnya
        sekaligus. Satu benda boleh berada pada beberapa wadah.
  \item Tiga perintah pemeriksaan tabrakan: \code{colliderect} untuk dua
        persegi, \code{spritecollide} untuk satu benda terhadap wadah, dan
        \code{groupcollide} untuk dua wadah.
  \item Mendeteksi tabrakan dan menyelesaikannya adalah dua pekerjaan yang
        berbeda. Pustaka menjawab yang pertama; aturan permainan Anda
        menentukan yang kedua.
  \item Penyelesaian tabrakan dilakukan sumbu demi sumbu: gerak mendatar
        diperiksa dan diselesaikan sampai tuntas, baru gerak tegak. Dengan
        begitu sisi tabrakan dapat dipastikan dari arah gerak.
  \item Posisi pada \code{rect} diperbaiki, lalu nilainya disalin kembali ke
        posisi pecahan, sehingga posisi tetap menjadi sumber kebenaran.
  \item Kecepatan pada sumbu yang diselesaikan dinolkan, agar benda tidak
        tertekan ke dalam bongkahan dan bergetar di permukaannya.
  \item Persegi pembatas kurang teliti untuk benda miring atau berlubang.
        \istilah{Mask} dipakai untuk keperluan itu, dan cara yang paling
        hemat adalah memakai persegi pembatas sebagai penyaring lebih dahulu.
  \item Penanganan tabrakan yang salah tidak menimbulkan pesan kesalahan.
        Empat gejala yang paling sering adalah benda menembus bongkahan tipis,
        benda terangkat ke atas bongkahan padahal menabrak sisinya, benda
        bergetar di permukaan, dan benda dinilai bersentuhan padahal belum.
  \item Cara menelusurinya: dapatkan gejala yang dapat diulang, persempit
        wilayahnya, pasang alat ukur, lalu ubah satu hal pada satu waktu.
  \item Pemeriksa mandiri di dalam program, misalnya penghitung tumpang
        tindih yang harus selalu bernilai nol, menemukan kesalahan lebih cepat
        daripada pengamatan mata.
  \item Bab ini menutup \CPMK{1}. Seluruh perkakas Blok~I---lingkungan kerja,
        siklus permainan, gerak, fisika, dan tabrakan---telah lengkap.
\end{itemize}
\end{rangkumanbox}
